
%\needspace{4\baselineskip}
\color{uniblue}{\subsection{Защита минимального напряжения (ЗМН)}}
\color{black}

\begin{enumerate}

\item
ЗМН отслеживает снижение напряжения на секции шин и формирует команду на отключение выключателя. В устройстве защита минимального напряжения представлена независимыми функциями (ступенями) «ЗМН 1» и «ЗМН 2».

\item
Функции могут применяться:
\begin{itemize}
\item в качестве индивидуальной/групповой защиты при потере питания секции шин, действующей на отключение присоединения;
\item для обнаружения провалов напряжения в ЛЭП и в электрических машинах.
\end{itemize}

\item
При снижении ИО напряжения ниже уставки активируются соответствующие выходные сигналы срабатывания ИО:
\begin{itemize}
    \item «ЗМН 1: ИО Uab»;
    \item «ЗМН 1: ИО Ubc»;
    \item «ЗМН 1: ИО Uca».
\end{itemize}

\item
Посредством уставки SGF2 предоставляется возможность настройки работы функции в одном из трех режимов:
\begin{itemize}
    \item  по снижению одного линейного напряжения;
    \item по снижению двух линейных напряжений;
    \item  по снижению трех линейных напряжений.
\end{itemize}

\item
Обязательным условием формирования выходных сигналов пуска и срабатывания функции является наличие активного сигнала «В: Включен», что исключает работу при отключенном выключателе и отсутствие активного сигнала блокировки от неисправности цепей напряжения ТН «БНН: Срабатывание».

\item
При наличии неисправности в цепях напряжения все функции автоматически блокируются.

\item
Уставкой SGF3 в функции предусмотрен контроль симметрии по напряжению обратной последовательности. В режиме работы по трем линейным напряжениям и увеличению напряжения ОП, пуск и срабатывание функции будет заблокировано.

\item
При наличии пусковых условий, срабатывание функции происходит после отсчета выдержки времени таймера T1.

\item
Ввод функций в работу на этапе параметрирования устройства осуществляется своим программным переключателем SGF1 «Ввод\_функции». Информация о введенной функции отображается выходным сигналом «ЗМН1: Ввод», «ЗМН2: Ввод».

\item
Функции могут быть выведены из работы оперативно путем активации сигнала «ОВ ЗМН 1» «ОВ ЗМН 2» соответственно, что характеризуется активным состоянием сигнала «ЗМН 1: Оперативный вывод», «ЗМН 2: Оперативный вывод».

\item
При активном состоянии входного сигнала «ЗМН 1 на сигн», «ЗМН 2 на сигн» действие соответствующей функции переводится на сигнал.

\item
Алгоритм работы функции «ЗМН 1» выполнен в соответствии с рисунком \ref{zmn:img1}. Алгоритм работы функции «ЗМН 2» аналогичный. В таблице \ref{zmn:tbl1} приведены параметры, необходимые для настройки функции ЗМН 1. В таблице \ref{zmn1:tbl1} приведены параметры, необходимые для настройки функции ЗМН 2.

\vspace{1mm}
\begin{figure}[H]
\centering
\includegraphics[width=1\textwidth,height=1\textheight,keepaspectratio]{\fbpath/112. ЗМН/_latex/img/zmn.pdf}
\captionof{figure}{Функциональная схема логики ЗМН 1}\label{zmn:img1}
\end{figure}


\end{enumerate}
